\begin{problem}[Matrix Recurrence and Contraction Mapping]
Define $a_1=0$ and, for $n\ge 2$, define the sequence by the scalar recurrence
\[
a_n=\frac{a_{n-1}+2}{a_{n-1}+3}.
\]


\begin{enumerate}
    \item Prove by induction that $0\le a_n<\sqrt{3}-1$ for all $n\ge1$.
    \item Let $L=\sqrt{3}-1$. Prove the convergence rate
    \[ |a_n-L|<\left(\tfrac{1}{3}\right)^{n-1}|a_1-L|. \]
    \item Determine whether the series $\sum_{n=1}^\infty (L-a_n)$ converges.
\end{enumerate}
\end{problem}

\begin{hint}
Show $L$ is the positive fixed point of $f(x)=\tfrac{x+2}{x+3}$ and note $f'(x)=\tfrac{1}{(x+3)^2}>0$. 
Express $a_n-L$ in terms of $a_{n-1}-L$ to find the contraction constant.
\end{hint}

\begin{solution}
We will use the explicit recurrence
\[
a_n=\frac{a_{n-1}+2}{a_{n-1}+3}.
\]
For (1) base case $a_1=0$ lies in the interval. If $0\le a_k< L$ then applying increasing $f$ gives $f(0)\le a_{k+1}<f(L)=L$, so $0\le a_{k+1}<L$. 

For (2) compute
\[
a_n-L=\frac{a_{n-1}+2}{a_{n-1}+3}-\frac{L+2}{L+3}=(a_{n-1}-L)\cdot\frac{1}{(a_{n-1}+3)(L+3)}.
\]
Because $a_{n-1}\ge0$ and $L+3>3$, the multiplier is $<\tfrac{1}{3}$, so iterating yields the claimed geometric bound.

For (3) set $u_n=L-a_n$. From (2) $|u_n|\le C(\tfrac{1}{3})^{n-1}$ for some constant $C$, so $\sum u_n$ is dominated by a convergent geometric series; hence the series converges absolutely.
\end{solution}

\begin{takeaways}
    \item A fixed point of a recurrence can be found by solving $x=f(x)$ (or the sequence: $x=\frac{x+2}{x+3}$). The behavior of the derivative $f'(x)$ determines monotonicity and depends on the fixed point(s).
\item Fixed-point iteration and derivative bounds give convergence rates.
\item Comparison with geometric series proves absolute convergence of the error series.
\item Fixed Points are Attractors: In sequences $a_n = f(a_{n-1})$, if $|f'(L)| < 1$, the sequence converges to $L$ (You can prove this)
\item Do you see the connection between this problem and eigenvalues of the matrix? 
\end{takeaways}
